% Standalone source-record documentation; no research attempts.
\documentclass[11pt]{article}
\usepackage[margin=25mm]{geometry}
\usepackage{fontspec,unicode-math}
\setmainfont{Latin Modern Roman}
\setmathfont{Latin Modern Math}
\usepackage{amsmath,amssymb,mathtools,xurl,hyperref}
\hypersetup{colorlinks=true,urlcolor=blue,linkcolor=blue}
\setcounter{secnumdepth}{0}
\setlength{\parindent}{0pt}
\setlength{\parskip}{5pt}
\setlength{\emergencystretch}{3em}
\newcommand{\sref}[2]{\hyperlink{source:#1:#2}{[S#2]}}
\newcommand{\cataloguescope}[1]{\paragraph{Source record.}#1}
\newcommand{\floor}[1]{\left\lfloor#1\right\rfloor}
\begin{document}
\section{Ensemble Necessity in Data Assimilation}\label{statement}

\subsection{Definitions and mathematical statement}

\paragraph{Shared notation.}Standard number systems are $\mathbb N,\mathbb Z,\mathbb Q,\mathbb R,\mathbb C$. Specialized notation and index ranges follow the cited source and the definitions below.


Data assimilation estimates a dynamical state by combining a forecast with observations. An ensemble method propagates several candidate states; a filtering error measures the accuracy of the resulting state estimate.
\paragraph{Record scope and qualifications.}
The catalogue title is retained as metadata. Named certificates, contracts, bounds or frameworks in that title are not asserted as results of the original source extract.
The supplied registry flags this extraction as unrecoverable. It is displayed as a literal source fragment; the definitions provide context rather than a replacement formal statement.

The experiment is evidence in a stated low-dimensional setting, not a universal no-ensemble guarantee; the excerpt is incomplete.

Source-record limitation: Source flags the statement as unrecoverable: The official report poses this question as a talk title and then gives a negative experimental answer in its studied setting: a learned single-state forecast attains ensemble-filter accuracy.
The source text below is rendered literally where necessary; damaged formulas and composite extracts are not silently repaired into a new mathematical claim.
The supplied catalogue marks this record as unrecoverable. Its extract and terminology are documented, but a complete original problem has not been reconstructed.
\paragraph{Statement recorded in the supplied source.}
Are ensemble-based data assimilation methods really necessary for accurate filtering? We investigate the ability to discover data assimilation (DA) schemes meant for chaotic dynamics with deep learning. The focus is on learning the analysis step of sequential DA, from state trajectories and their observations, using a simple residual convolutional neural network, while assuming the dynamics to be known. Experiments are performed with low-order dynamics which display spatiotemporal chaos and for which solid benchmarks for DA performance exist. The accuracy of the states obtained from the learned analysis approaches that of the best possibly tuned ensemble Kalman filter, and is far better than that of variational DA alter- natives. Critically, this can be achieved while propagating even just a single state in the forecast step. We investigate the reason for achieving ensemble filtering accuracy without an ensemble. We diagnose that the analysis scheme actually identifies key dynamical perturbations, mildly aligned with the unstable subspace, from the forecast state alone, without any ensemble-based covariances representa- tion. This reveals that the analysis scheme has learned some multiplicative ergodic \sref{statement}{2}
\paragraph{Source status.}
The supplied catalogue labels this source record partially solved; that metadata was checked against the supplied snapshot on 2026-10-11. The label is preserved for provenance and is not an independent verification of present open status or mathematical correctness. The catalogue literature assessment is dated 2026-08-22.
\subsection{Short English statement}
The fragment questions whether ensembles are necessary and reports a learned single-state filtering experiment.
\subsection{Sources}
\begin{itemize}
\item[\hypertarget{source:statement:1}{S1}] Ulam AI and the UnsolvedMath Contributors, UnsolvedMath: A Curated Collection of Open Mathematics Problems, catalogue entry. Catalogue curation and metadata: CC BY 4.0. The source extract is retained; definitions, English summaries and qualifications are editorial additions.
\url{https://huggingface.co/datasets/ulamai/UnsolvedMath}.
\item[\hypertarget{source:statement:2}{S2}] Data Assimilation: From Mathematical and Statistical Foundations to Applications (2025), Oberwolfach Reports, DOI 10.4171/owr/2025/10, pp. 451--484
\url{https://doi.org/10.4171/owr/2025/10}.
\item[\hypertarget{source:statement:3}{S3}] Bocquet et al., ISDA 2025 abstract
\url{https://www.data-assimilation.riken.jp/en/events/isda_2021/abstracts/isda2025_03.pdf}.
\item Source-specific Oberwolfach Reports licences, including CC BY-SA 4.0 where applicable, are retained; see THIRD\_PARTY\_NOTICES.md.
\end{itemize}
\paragraph{Preparation and review.}
The original source and UnsolvedMath contributors retain attribution for the source extract. Codex (GPT-6; exact serving revision not exposed) assisted with definitions, English summaries, scope notes and formatting. No new human mathematical authorship or independently verified solution is asserted. Catalogue identity checked against the complete supplied registry and current repository. Distinct upstream records, including repeated questions and broad or damaged extracts, are retained for complete coverage of the source CS category. A repeated record is not a new independent conjecture. Complete mathematical-equivalence and current-literature review has not been performed. Preparation and final review follow contribution prompts C and D, with source gaps disclosed. Structural and rendering checks do not establish a new theorem or open status.
\label{end:statement}
\end{document}
